% §8 Exercises
%
% Cross-reference resolution status:
%
% DEFINED here:
%   \label{ex:ch8-1} through \label{ex:ch8-11}
%
% RESOLVED (used here, defined elsewhere):
%   \Cref{sec:entropy-refresher}                  -- §8.1 (Plan 8 Task 2)
%   \Cref{eq:binary-entropy}                      -- §8.1 (Plan 8 Task 2)
%   \Cref{eq:mutual-info}                         -- §8.1 (Plan 8 Task 2)
%   \Cref{sec:bsc-capacity}                       -- §8.2 (Plan 8 Task 3)
%   \Cref{prop:bsc-capacity}                      -- §8.2 (Plan 8 Task 3)
%   \Cref{sec:entropy-floor}                      -- §8.3 (Plan 8 Task 4)
%   \Cref{thm:entropy-floor}                      -- §8.3 (Plan 8 Task 4)
%   \Cref{sec:bloom-gap}                          -- §8.4 (Plan 8 Task 5)
%   \Cref{prop:bloom-overhead}                    -- §8.4 (Plan 8 Task 5)
%   \Cref{def:union-of-sets}                      -- §6.1 (Plan 6)
%   \Cref{def:bsc}                                -- §2.1 (Plan 2)
%   \Cref{thm:bloom-optimal-k}                    -- §5.2 (Plan 5)
%
% Difficulty-marker macros (not in alex.sty; defined locally if not yet defined):
\providecommand{\exstar}{$\star$}
\providecommand{\exstarstar}{$\star\star$}
\providecommand{\exstarstarstar}{$\star\star\star$}

\section*{Exercises}
\addcontentsline{toc}{section}{Exercises}

\begin{enumerate}

% -----------------------------------------------------------------------
% WARM-UP  (★ × 3): exercises 1, 2, 3
% -----------------------------------------------------------------------

\item \exstar
\label{ex:ch8-1}
\textbf{Entropy of a biased coin.}
Compute the binary entropy $H(p)$ from \Cref{eq:binary-entropy} for
$p \in \{0.01, 0.5, 0.99\}$. Report each to four decimal places. In
one sentence, name the symmetry $H(p) = H(1 - p)$ that makes the
$p = 0.01$ and $p = 0.99$ values agree, and identify the maximum of
$H$ on $[0, 1]$.

\item \exstar
\label{ex:ch8-2}
\textbf{BSC capacity computation.}
For $\fprate \in \{0.01, 0.1, 0.5\}$, compute the BSC channel
capacity $C = 1 - H(\fprate)$ from \Cref{prop:bsc-capacity}. Report
each to four decimal places. In one sentence, identify the regime
($\fprate$ near $0$ vs $\fprate$ near $1/2$) where the channel
carries close to a full bit per use vs where it carries essentially
nothing.

\item \exstar
\label{ex:ch8-3}
\textbf{Entropy floor evaluation.}
For each target false-positive rate $\fprate \in \{0.1, 0.01,
0.001\}$, compute the entropy floor $-\log_2 \fprate$ in bits per
element from \Cref{thm:entropy-floor}. Report each to four decimal
places. In one sentence, describe the linear-in-$-\log_2 \fprate$
scaling of the floor and contrast it with the linear-in-$1/\fprate$
scaling one might naively expect from the rate alone.

% -----------------------------------------------------------------------
% MEDIUM  (★★ × 5): exercises 4, 5, 6, 7, 8
% -----------------------------------------------------------------------

\item \exstarstar
\label{ex:ch8-4}
\textbf{Notebook-tied gap plot.}
Open the chapter 8 companion notebook
\texttt{python/notebooks/ch08-information-bounds.ipynb}.
\begin{enumerate}
  \item[(a)] For $\fprate \in [10^{-5}, 10^{-1}]$ (log-spaced, at
    least 25 points), compute the entropy floor $-\log_2 \fprate$
    bits per element and the standard Bloom cost
    $-\log_2 \fprate \cdot \log_2 e$ bits per element from
    \Cref{prop:bloom-overhead}.
  \item[(b)] Plot both as line plots on a log-$\fprate$ axis.
    Visually confirm that the two lines are parallel on this scale
    (they differ by a constant additive offset in bits per
    element).
  \item[(c)] Plot the ratio (Bloom cost) / (floor) vs $\fprate$.
    Verify the ratio is constant at $\approx 1.443$ across the
    entire range, matching $\log_2 e = 1/\ln 2$ exactly. In one
    sentence, state which feature of the standard Bloom filter
    (the $k$ independent hash functions of \Cref{thm:bloom-optimal-k})
    produces this constant overhead.
\end{enumerate}

\item \exstarstar
\label{ex:ch8-5}
\textbf{Mutual information of the BSC.}
The BSC of \Cref{def:bsc} has input $X \in \{\top, \bot\}$ with
distribution $P_X$, output $Y \in \{\top, \bot\}$, and crossover
$\fprate$. Use \Cref{eq:mutual-info} as $I(X; Y) = H(Y) - H(Y |
X)$.
\begin{enumerate}
  \item[(a)] Show that $H(Y | X) = H(\fprate)$ for every $P_X$ (the
    conditional output uncertainty is just the channel's flip
    uncertainty, independent of the input distribution).
  \item[(b)] Compute the marginal output distribution as $P_Y(\top)
    = P_X(\top)(1 - \fprate) + P_X(\bot) \fprate$, and substitute
    to obtain $I(X; Y) = H(P_X(\top)(1 - \fprate) + P_X(\bot)
    \fprate) - H(\fprate)$.
  \item[(c)] Maximize $I(X; Y)$ over $P_X$ by maximizing $H(Y)$ in
    the first term. Since $H$ on a binary alphabet is maximized at
    $1/2$, the optimum is $P_X$ uniform, which produces $P_Y$
    uniform and $H(Y) = 1$. Recover the BSC capacity
    $C = 1 - H(\fprate)$ of \Cref{prop:bsc-capacity}.
\end{enumerate}

\item \exstarstar
\label{ex:ch8-6}
\textbf{The 1.44 constant.}
The constant $\log_2 e \approx 1.4427$ appears in
\Cref{prop:bloom-overhead} as the ratio of standard Bloom cost to
entropy floor.
\begin{enumerate}
  \item[(a)] Starting from the change-of-base identity $\log_2 e =
    \ln e / \ln 2 = 1 / \ln 2$, compute $\log_2 e$ to four decimal
    places using $\ln 2 \approx 0.6931$.
  \item[(b)] Verify the equivalent form $\log_2(1/\fprate) \cdot
    \log_2 e = \ln(1/\fprate) / (\ln 2)^2$ by substituting
    $\log_2 x = \ln x / \ln 2$ on both factors. State why this
    quadratic-in-$1/\ln 2$ form is the natural one for the Bloom
    filter's bits-per-element formula (one factor of $1/\ln 2$
    comes from the optimal-$k$ scaling, a second from the
    conversion of natural to binary logarithms).
  \item[(c)] In one sentence, state the geometric meaning of
    $\log_2 e$: it is the number of base-$2$ digits required to
    represent one nat of information.
\end{enumerate}

\item \exstarstar
\label{ex:ch8-7}
\textbf{Entropy floor for sparse universes.}
\Cref{thm:entropy-floor} is stated in the asymptotic regime $N
\gg n$, where $N = |\Universe|$ is the universe size and $n = |A|$
is the set size.
\begin{enumerate}
  \item[(a)] Suppose instead that $N$ is comparable to $n$, say
    $N = 10 n$. Recompute the counting argument of
    \Cref{sec:entropy-floor}: how many false positives can a
    representation accept before the FPR ceiling $\fprate$ is
    breached? Does the dominant term in the bound $|A_R| \le n +
    \fprate (N - n)$ change as $N$ approaches $n$?
  \item[(b)] Discuss the regime where the asymptotic form
    $-n \log_2 \fprate$ is accurate (large $N$, small $\fprate$)
    vs where finite-$N$ corrections matter (small $N/n$ ratio, or
    $\fprate$ close to the trivial ceiling $1 - n/N$).
  \item[(c)] In one sentence, identify the trivial representation
    that achieves $\fprate = 1 - n/N$ at zero space cost (always
    return $\top$), and argue that the entropy floor is
    informative only below this trivial ceiling.
\end{enumerate}

\item \exstarstar
\label{ex:ch8-8}
\textbf{Two-sided floor.}
\Cref{thm:entropy-floor} is stated for FNR-zero approximate sets.
A two-sided approximate set permits both false positives at rate
$\fprate$ and false negatives at rate $\fnrate$.
\begin{enumerate}
  \item[(a)] Sketch (without proof) why a two-sided set can have a
    lower space cost than $-n \log_2 \fprate$ at the same maximum
    error. Hint: the two-sided set splits its error budget across
    both directions, so each side can be allowed a larger rate
    than the one-sided $\fprate$.
  \item[(b)] Suppose the maximum-error budget is $\fprate = \delta$
    for the one-sided case and $\max(\fprate, \fnrate) = \delta$
    for the two-sided case. Argue informally that the two-sided
    set need only achieve $\fprate = \fnrate = \delta$ (each
    saturating its own ceiling), so the relevant rate in the
    bound is now $\delta$ on each side rather than $\delta$
    concentrated on the FPR axis alone.
  \item[(c)] In one sentence, state the qualitative consequence: a
    two-sided set with $\fprate = \fnrate$ can sit below the
    one-sided floor at the same max-error budget, at the cost of
    abandoning the FNR-zero guarantee of \Cref{thm:entropy-floor}.
\end{enumerate}

% -----------------------------------------------------------------------
% CHALLENGING / EXTENSION  (★★★ × 3): exercises 9, 10, 11
% -----------------------------------------------------------------------

\item \exstarstarstar
\label{ex:ch8-9}
\textbf{The floor as a counting argument.}
\Cref{thm:entropy-floor} is proved in \Cref{sec:entropy-floor} via
a counting argument.
\begin{enumerate}
  \item[(a)] Reproduce the proof sketch in your own words. Begin
    with: a representation $R$ of length $L$ bits induces a query
    algorithm that accepts a superset $A_R \supseteq A$; FNR-zero
    forces the inclusion; FPR $\le \fprate$ caps $|A_R \setminus
    A|$.
  \item[(b)] Identify the load-bearing step: the pigeonhole bound
    on the number of distinct $A$ that a single $R$ can encode.
    State the inequality $2^L \cdot (\text{accepting sets per
    } R) \ge \binom{N}{n}$ and explain in one sentence why each
    of the $\binom{N}{n}$ possible truth-sets $A$ must be encoded
    by at least one $R$.
  \item[(c)] Take the binary logarithm of the resulting inequality
    and apply Stirling's approximation to $\binom{N}{n}$ to
    recover $L \ge n \log_2(1/\fprate)$ asymptotically. Identify
    which step of the derivation discards the additive
    lower-order terms.
\end{enumerate}

\item \exstarstarstar
\label{ex:ch8-10}
\textbf{Compressed Bloom filter.}
\cite{mitzenmacher-2002} showed that arithmetic compression of a
standard Bloom filter's bit array can recover much of the $\log_2
e$ overhead identified in \Cref{prop:bloom-overhead}.
\begin{enumerate}
  \item[(a)] Sketch the argument: under the optimal $k$ of
    \Cref{thm:bloom-optimal-k}, the bitmap is tuned so that
    approximately half its bits are $1$ and half are $0$ at
    insertion saturation. Conclude that the per-bit entropy of
    the bitmap is $H(1/2) = 1$ bit, so the bitmap is already
    informationally dense per stored bit.
  \item[(b)] Argue that the $\log_2 e$ overhead comes from a
    different place: not from per-bit redundancy in the bitmap,
    but from the choice of bitmap-with-$k$-independent-hashes as
    the underlying representation. Arithmetic compression
    eliminates the per-bit redundancy that DOES exist when the
    bitmap is loaded below the half-full optimum (e.g., during
    construction or with a non-optimal $k$).
  \item[(c)] In one sentence, identify why full elimination of
    the gap still requires abandoning the bitmap-with-$k$-hashes
    structure: the standard Bloom representation is wasteful at
    the structural level, not merely at the per-bit-encoding
    level, and only a different representation (chapter 14's BHF)
    can close the gap to the floor of \Cref{thm:entropy-floor}.
\end{enumerate}

\item \exstarstarstar
\label{ex:ch8-11}
\textbf{Entropy floor for the union.}
\Cref{def:union-of-sets} of §6.1 derived the FPR of the union
$\tilde A \cup \tilde B$ of two approximate sets. The §8.3 floor
applies to the union itself as an FNR-zero approximate set.
\begin{enumerate}
  \item[(a)] Apply \Cref{thm:entropy-floor} to the union: the bits
    required to represent $A \cup B$ approximately at FPR
    $\fprate_\cup$ are at least $|A \cup B| \cdot \log_2(1 /
    \fprate_\cup)$.
  \item[(b)] Compare to the cost of representing $A$ and $B$
    separately at FPRs $\fprate_A, \fprate_B$ and combining their
    query algorithms by disjunction. The separate cost is at
    least $|A| \log_2(1/\fprate_A) + |B| \log_2(1/\fprate_B)$ bits;
    the combined query has FPR $\le \fprate_A + \fprate_B -
    \fprate_A \fprate_B$ (and equals $\fprate_A + \fprate_B$ when
    both are small). When $A$ and $B$ are disjoint and
    $\fprate_A = \fprate_B = \fprate_\cup / 2$ (so the combined
    rate matches the union's $\fprate_\cup$), show that the two
    costs agree to leading order: $(|A| + |B|)
    \log_2(2/\fprate_\cup) = |A \cup B|(\log_2(1/\fprate_\cup) +
    1)$, exceeding the union floor by exactly one bit per element.
  \item[(c)] Identify when union representation is strictly more
    efficient: when $A$ and $B$ overlap heavily ($|A \cup B| \ll
    |A| + |B|$), the union floor scales with $|A \cup B|$, while
    the separate-then-combine cost scales with $|A| + |B|$. State
    in one sentence the practical guideline: store the union
    directly when membership in $A$ vs $B$ is not separately
    required by downstream queries.
\end{enumerate}

\end{enumerate}
